\scanpage{11}
\textbf{C. R. Acad. Sc. Paris, t. 279 (16 septembre 1974)}\hfill\textbf{Série A — 447}

\rule{\linewidth}{0.4pt}

ALGÈBRE. — \emph{Théorème de Skolem-Nother pour les algèbres d'Azumaya
graduées sur un anneau semi-local.} Note (*) de M. \textbf{Olivier Laborde},
présentée par M. René Garnier.

Le théorème de Skolem-Nother se généralise au cas des algèbres d'Azumaya graduées grâce
à quelques propriétés des anneaux gradués simples et semi-simples. Une première application
à la K-théorie de ce résultat est la régularité de ces algèbres pour le produit tensoriel.

1. \textsc{Modules et anneaux gradués semi-simples}. — Un anneau, A, est gradué si additivement,
il est une somme directe de sous-groupes $\mathrm A=\mathrm A_0\oplus\mathrm A_1$ qui, pour la multiplication
vérifient les inclusions $\mathrm A_i.\mathrm A_j\subseteq\mathrm A_{i+j}$ ($i,j\in\mathbf Z/2\mathbf Z$). On notera $|\mathrm A|$ l'anneau sous-jacent
non gradué.

Un A-module gradué, M, est un $|\mathrm A|$-module qui est pour l'addition une somme directe
$\mathrm M=\mathrm M_0\oplus\mathrm M_1$ qui satisfait $\mathrm A_i\,\mathrm M_j\subseteq\mathrm M_{i+j}$ ($i,j\in\mathbf Z/2\mathbf Z$) et $|\mathrm M|$ sera le $|\mathrm A|$-module
sous-jacent non gradué.

Dans ce qui suit les anneaux, les idéaux, les algèbres et les modules seront, sauf mention
expresse du contraire, gradués et on les appellera des Anneaux, des Idéaux, des Algèbres
et des Modules. Pour d'autres renseignements, \emph{voir} (${}^3$).

\textsc{Définition 1}. — Un A-Module à gauche M est Simple si les seuls sous-Modules de M
sont $\{0\}$ et M.

\textsc{Proposition 1}. — \emph{Un A-module à gauche Simple non nul est de la forme
$\mathrm A\,\unsure{s/}\mathbf m$ où $\mathbf m$
est un Idéal à gauche maximal de} A.

\emph{Preuve}. — Il est immédiat que $\mathrm A\,\unsure{s/}\mathbf m$ est Simple. Réciproquement, soit M un A-Module
Simple et $x$ un élément homogène non nul de M. Considérons l'application $f$ qui à $\alpha\in\mathrm A$
fait correspondre $\alpha x\in\mathrm M$. Le noyau de $f$ est un Idéal de A quel que soit le degré de $x$ :
Si $\alpha\in\operatorname{Ker}(f)$, on a
\[
\alpha x=\alpha_0x+\alpha_1x=0\quad\Rightarrow\quad\alpha_0x=\alpha_1x=0
\qquad\text{et}\qquad\alpha_0\text{ et }\alpha_1
\]
appartiennent eux aussi à $\operatorname{Ker}(f)$.

\textsc{Lemme}. — \emph{Soit} M \emph{un A-Module qui est somme d'une famille}
$(\mathrm M_i)_{i\in\unsure{i|I}}$ \emph{de sous-Modules
Simples et} F \emph{un sous-Module de} M. \emph{Alors il existe}
$\mathrm J\subseteq\mathrm I$ \emph{tel que}
$\mathrm M=\mathrm F\oplus\bigoplus_{j\in\mathrm J}\mathrm N_j$.

La démonstration est identique à celle du cas non gradué.

\textsc{Proposition 2}. — \emph{Soit} M \emph{un A-Module. Les propositions suivantes sont équivalentes} :

\emph{(a)} M \emph{est somme d'une famille de sous-Modules Simples.}

\emph{(b)} M \emph{est somme directe d'une famille de sous-Modules Simples.}

\emph{(c) Tout sous-Module de} M \emph{est facteur direct.}

\emph{Preuve}. — Elle est rendue immédiate par le lemme précédent.

\textsc{Définition 2}. — On dira qu'un A-module M est Semi-Simple s'il vérifie l'une des conditions
équivalentes de la proposition 2.

\mspage{10}
